\documentclass[11pt,a4paper]{article}
\usepackage{fontspec,amsmath,amssymb,amsthm,microtype}
\setmainfont{lmroman10-regular.otf}[BoldFont=lmroman10-bold.otf,ItalicFont=lmroman10-italic.otf,BoldItalicFont=lmroman10-bolditalic.otf]
\usepackage{xcolor}
\usepackage[most]{tcolorbox}
\usepackage{mathrsfs}
\providecommand{\lt}{<}
\providecommand{\gt}{>}
\setmonofont{DejaVuSansMono.ttf}[BoldFont=DejaVuSansMono-Bold.ttf,ItalicFont=DejaVuSansMono-Oblique.ttf,BoldItalicFont=DejaVuSansMono-BoldOblique.ttf,Scale=0.86]
\newfontfamily\leanfallback{STIXTwoMath-Regular.otf}[Scale=0.86]
\definecolor{leanrule}{gray}{0.62}
\definecolor{leanback}{gray}{0.962}
\definecolor{leantitle}{gray}{0.22}
\definecolor{leanmuted}{gray}{0.38}
\newtcolorbox{leanbox}[1]{enhanced,breakable,colback=leanback,colframe=leanrule,
  boxrule=0.4pt,arc=0.6mm,left=2.4mm,right=2.4mm,top=1.8mm,bottom=1.8mm,
  before skip=0.8em,after skip=1em,fonttitle=\small\bfseries,coltitle=leantitle,
  colbacktitle=leanback,titlerule=0.3pt,title={#1}}
\newlength{\leancharwidth}
\newenvironment{leancode}{\par\smallskip\begingroup\ttfamily\footnotesize\raggedright
  \settowidth{\leancharwidth}{M}\setlength{\parskip}{0pt}\setlength{\parindent}{0pt}
  \frenchspacing}{\par\endgroup\smallskip}
\newcommand{\leanline}[2]{\hangindent=\dimexpr#1\leancharwidth+4\leancharwidth\relax
  \hangafter=1\noindent\hspace*{#1\leancharwidth}#2\par}
\newcommand{\leansource}[1]{\par{\footnotesize\color{leanmuted}#1}\par}
\newcommand{\leanlabel}[1]{\par\smallskip\noindent{\small\bfseries #1}\par}
\usepackage[margin=25mm]{geometry}
\usepackage[unicode,hidelinks]{hyperref}
\hypersetup{pdfauthor={LangevinLean Project}}
\setlength{\emergencystretch}{3em}
\setlength{\parskip}{0.4em}
\setlength{\parindent}{0pt}
\allowdisplaybreaks[2]
\raggedbottom
\title{\texorpdfstring{Weighted Dynkin identity for the constructed Brownian Langevin solution}{Weighted Dynkin identity for the constructed Brownian Langevin solution}}
\author{LangevinLean Project}
\date{}
\begin{document}\maketitle
Draft mathematical exposition; these notes may contain mistakes.


This module completes the stochastic bridge for compact \(C^{2}\) tests. All formulas are derived from the actual Brownian innovations, the continuous \(L^{2}\) integral solution, and finite Gaussian Euler chains.



The actual Brownian Euler chains have an eventual uniform score bound on every time partition:


\[\exists R\ge0\ \exists N\ \forall n\ge N\ \forall k\le n+1:
\quad\|\nabla f(E_k^{(n)})\|_{L^2(P)}\le R.\]


Actual Gaussian weak consistency and the mesh limit prove, for bounded initial weights independent of the innovations and \(\phi\in C_c^2\),


\[\mathbb E[w\phi(Y_T)]-\mathbb E[w\phi(Y_0)]
=\int_0^T\mathbb E[w\mathcal L_f\phi(Y_t)]\,dt.\]

Earlier results used here: \cite[Section~3.42]{LangevinLean2026ULASC}.

\section{\texorpdfstring{Bound all actual Brownian Euler scores}{Bound all actual Brownian Euler scores}}


Use the deterministic uniform grid bound in \(L^{2}\), with the lifted gradient as drift and the actual Brownian noise path. The established equality between additive matching-noise Euler states and the Euler chain in the Hilbert space of square-integrable random vectors identifies the objects used in the probabilistic weak-step proof.

\begin{leanbox}{Lean: \texttt{brownian\_\allowbreak{}euler\_\allowbreak{}score\_\allowbreak{}grid\_\allowbreak{}eventually\_\allowbreak{}bounded}}
\leansource{\texttt{brownian\_\allowbreak{}euler\_\allowbreak{}score\_\allowbreak{}grid\_\allowbreak{}eventually\_\allowbreak{}bounded} \quad AverageLMC/\allowbreak{}LaplacianBrownianDynkin.lean, lines 28–48}
\begin{leancode}
\leanline{0}{theorem\ brownian\_euler\_score\_grid\_eventually\_bounded}
\leanline{4}{(f\ :\ EuclideanSpace\ ℝ\ ι\ →\ ℝ)\ (L\ :\ ℝ≥0)\ (hLip\ :\ LipschitzWith\ L\ (gradient\ f))}
\leanline{4}{(initial\ :\ Lp\ (EuclideanSpace\ ℝ\ ι)\ 2\ P)}
\leanline{4}{(B\ :\ ι\ →\ ℝ≥0\ →\ Ω\ →\ ℝ)\ (hB\ :\ ∀\ i,\ IsPreBrownianReal\ (B\ i)\ P)}
\leanline{4}{(Y\ :\ ℝ\ →\ Lp\ (EuclideanSpace\ ℝ\ ι)\ 2\ P)\ (hY\ :\ Continuous\ Y)}
\leanline{4}{(T\ :\ ℝ)\ (hT\ :\ 0\ <\ T)}
\leanline{4}{(hYe\ :\ ∀\ t\ ∈\ Icc\ 0\ T,\ Y\ t\ -\ initial\ =}
\leanline{6}{-(∫\ s\ in\ 0..t,\ ULA.lpDrift\ P\ (gradient\ f)\ hLip\ (Y\ s))\ +}
\leanline{8}{ULA.brownianLangevinNoise\ P\ B\ hB\ 0\ t)\ :}
\leanline{4}{∃\ R\ :\ ℝ,\ 0\ ≤\ R\ ∧\ ∃\ N\ :\ ℕ,\ ∀\ n\ ≥\ N,\ ∀\ k\ ≤\ n\ +\ 1,}
\leanline{6}{‖ULA.lpDrift\ P\ (gradient\ f)\ hLip}
\leanline{8}{(ULA.Hilbert.lpChain\ P\ f\ (T\ /\ ((n\ +\ 1\ :\ ℕ)\ :\ ℝ))\ hLip\ initial}
\leanline{10}{(ULA.Hilbert.brownianEulerNoise\ P\ B\ hB\ (T\ /\ ((n\ +\ 1\ :\ ℕ)\ :\ ℝ)))\ k)‖\ ≤\ R}
\end{leancode}
\leanlabel{How the Lean proof follows the argument}
\texttt{brownian\_\allowbreak{}euler\_\allowbreak{}score\_\allowbreak{}grid\_\allowbreak{}eventually\_\allowbreak{}bounded} gives the formal counterpart of this step. Named proof applications include \texttt{h\allowbreak{}N}; the adjacent proof shows their hypotheses and order of use.
\end{leanbox}


Earlier results used here: \cite[Section~3.13.2]{LangevinLean2026ULASC}; \cite[Section~3.13.5]{LangevinLean2026ULASC}; \cite[Section~3.19.1]{LangevinLean2026ULASC}; \cite[Section~3.25.1]{LangevinLean2026ULASC}; \cite[Section~3.25.3]{LangevinLean2026ULASC}; \cite[Section~3.29.7]{LangevinLean2026ULASC}; \cite[Section~3.42.5]{LangevinLean2026ULASC}.

\section{\texorpdfstring{Derive weighted Dynkin without an assumed formula}{Derive weighted Dynkin without an assumed formula}}


For positive \(T\), actual Brownian coordinate increments give independent standard Gaussian vector innovations on every refining partition. Keep the bounded initial weight paired with the initialization, so each next innovation is fresh. Apply compact \(C^{2}\) Gaussian weak consistency, sum its residual using the uniform score moment, and pass to the limit with bounded uniformly continuous test and generator observations. The zero-time case follows directly from the integral equation. The conclusion is \(\mathbb E[w\varphi(Y_T)]-\mathbb E[w\varphi(Y_0)]=\int_0^T\mathbb E[w\mathcal L\varphi(Y_t)]\,dt\).

\begin{leanbox}{Lean: \texttt{lp\_\allowbreak{}langevin\_\allowbreak{}weighted\_\allowbreak{}dynkin\_\allowbreak{}compact}}
\leansource{\texttt{lp\_\allowbreak{}langevin\_\allowbreak{}weighted\_\allowbreak{}dynkin\_\allowbreak{}compact} \quad AverageLMC/\allowbreak{}LaplacianBrownianDynkin.lean, lines 54–110}
\begin{leancode}
\leanline{0}{theorem\ lp\_langevin\_weighted\_dynkin\_compact}
\leanline{4}{(f\ :\ EuclideanSpace\ ℝ\ ι\ →\ ℝ)\ (L\ :\ ℝ≥0)\ (hLip\ :\ LipschitzWith\ L\ (gradient\ f))}
\leanline{4}{(initial\ :\ Lp\ (EuclideanSpace\ ℝ\ ι)\ 2\ P)}
\leanline{4}{(B\ :\ ι\ →\ ℝ≥0\ →\ Ω\ →\ ℝ)\ (hB\ :\ ∀\ i,\ IsPreBrownianReal\ (B\ i)\ P)}
\leanline{4}{(w\ :\ Ω\ →\ ℝ)\ (hw\ :\ Measurable\ w)\ (A\ :\ ℝ)\ (hA\ :\ 0\ ≤\ A)}
\leanline{4}{(hwA\ :\ ∀ᵐ\ a\ ∂P,\ |w\ a|\ ≤\ A)}
\leanline{4}{(hscalar\ :\ ∀\ r\ :\ ℝ,\ 0\ <\ r\ →\ iIndepFun\ (fun\ q\ :\ ℕ\ ×\ ι\ ↦\ fun\ a\ ↦}
\leanline{6}{ULA.Hilbert.brownianEulerNoise\ P\ B\ hB\ r\ q.1\ a\ q.2)\ P)}
\leanline{4}{(hind\ :\ ∀\ r\ :\ ℝ,\ 0\ <\ r\ →\ IndepFun\ (fun\ a\ ↦\ (w\ a,\ initial\ a))}
\leanline{6}{(fun\ a\ k\ ↦\ ULA.Hilbert.brownianEulerNoise\ P\ B\ hB\ r\ k\ a)\ P)}
\leanline{4}{(Y\ :\ ℝ\ →\ Lp\ (EuclideanSpace\ ℝ\ ι)\ 2\ P)\ (hY\ :\ Continuous\ Y)}
\leanline{4}{(T\ :\ ℝ)\ (hT\ :\ 0\ ≤\ T)}
\leanline{4}{(hYe\ :\ ∀\ t\ ∈\ Icc\ 0\ T,\ Y\ t\ -\ initial\ =}
\leanline{6}{-(∫\ s\ in\ 0..t,\ ULA.lpDrift\ P\ (gradient\ f)\ hLip\ (Y\ s))\ +}
\leanline{8}{ULA.brownianLangevinNoise\ P\ B\ hB\ 0\ t)}
\leanline{4}{(φ\ :\ EuclideanSpace\ ℝ\ ι\ →\ ℝ)\ (hφ\ :\ ContDiff\ ℝ\ 2\ φ)\ (hc\ :\ HasCompactSupport\ φ)\ :}
\leanline{4}{lpWeightedTest\ P\ w\ φ\ (Y\ T)\ -\ lpWeightedTest\ P\ w\ φ\ initial\ =}
\leanline{6}{∫\ t\ in\ 0..T,\ lpWeightedTest\ P\ w\ (ULA.langevinGenerator\ f\ φ)\ (Y\ t)}
\end{leancode}
\leanlabel{How the Lean proof follows the argument}
\texttt{lp\_\allowbreak{}langevin\_\allowbreak{}weighted\_\allowbreak{}dynkin\_\allowbreak{}compact} separates the cases \texttt{h\allowbreak{}Tzero : T = 0} and its negation before combining them. The main intermediate facts are \texttt{h\allowbreak{}Tpos}, \texttt{htest}, \texttt{hgen}, \texttt{hmesh}, \texttt{h}, \texttt{hh} in that order. Named proof applications include \texttt{weighted\_\allowbreak{}dynkin\_\allowbreak{}of\_\allowbreak{}weak\_\allowbreak{}euler\_\allowbreak{}consistency}; the adjacent proof shows their hypotheses and order of use.
\end{leanbox}


Earlier results used here: \cite[Section~3.13.2]{LangevinLean2026ULASC}; \cite[Section~3.13.5]{LangevinLean2026ULASC}; \cite[Section~3.25.1]{LangevinLean2026ULASC}; \cite[Section~3.25.3]{LangevinLean2026ULASC}; \cite[Section~3.29.7]{LangevinLean2026ULASC}; \cite[Section~3.42.5]{LangevinLean2026ULASC}; \cite[Section~3.42.6]{LangevinLean2026ULASC}; \cite[Section~3.46.2]{LangevinLean2026ULASC}.

\begin{thebibliography}{ULA-LSI}
\bibitem[ULA-SC]{LangevinLean2026ULASC}
LangevinLean Project. \emph{Unadjusted Langevin under Strong Convexity}. LL-01, 2026. Draft report, content revision d50c0f64815b. \url{https://odakyildiz.github.io/LangevinLean/docs/pdf/UnadjustedLangevin/FullProof.pdf}.
\end{thebibliography}

\end{document}
